\subsection{\problemblock{*Constants} Data Block}\label{sec:block-constants}

Tables can be functions of user-defined variables as shown in
\cref{sec:table-user-defined-variables}. If user-defined variables have been referenced
in tables, values for those variables must be provided. They can be given in
\problemblock{*aero}, \problemblock{*prop}, \problemblock{*cg}, or
\problemblock{*constants} data blocks. Values are required only for variables used
by tables referenced in the current segment; variables belonging only to unused
tables need not be supplied.

The \problemblock{*constants} data block is optional. It is provided as a
convenience so that values for all user-defined variables can be grouped together.
The block is used only to provide such values. If the selected tables are not
functions of user-defined variables, a \problemblock{*constants} block should not
be input.

There are two styles for supplying these values. One style puts each value in the
data block where the associated table is referenced. For example, an
\problemblock{*aero} block that names a table containing a user-defined variable
also gives that variable's value. This style uses the \problemblock{*aero},
\problemblock{*prop}, and \problemblock{*cg} blocks. The other style places all
user-defined values in a single \problemblock{*constants} block. This groups them
in one place. Both styles are supported.

The keyword is followed by assignments. For example,

\begin{lstlisting}[style=taosmanual]
*constants type=3, tmult=1.2,
           mmult=1.15
\end{lstlisting}

sets \texttt{type} equal to 3.0, \texttt{tmult} equal to 1.2, and \texttt{mmult}
equal to 1.15. These variables must occur in the tables supplied with the problem.
Once a value has been given for a user-defined variable, it remains set until changed
in a subsequent data block. Values therefore need not be repeated in every segment
when they do not change; nevertheless, repeating them in each segment is good
practice because it makes the intended configuration explicit and helps avoid input
errors.
